

\noindent\textbf{LLM Agents.} LLM-based agents enhance reasoning by integrating external modules for planning, memory, and tool use. The ReAct \citep{yao2022react} architecture established the iterative ``reason-act'' paradigm, which has since evolved into specialized research on test-time searching \citep{yao2023tree, besta2024graph}, tool tuning \citep{zeng2025tool, qian2025toolrl, jiang2025verltool}, and memory structuring \citep{xu2025mem, chhikara2025mem0}. Beyond modular improvements, recent frameworks introduce diverse coordination strategies: AutoAgent \citep{chen2023autoagents} and OWL \citep{hu2025owl} utilize dynamic multi-agent coordination and role-playing for task decomposition, while Alita \citep{qiu2025alita} enables autonomous expansion via refined task-specific MCPs. Finally, Smolagents \citep{smolagents} prioritizes efficiency through lightweight, code-based actions and local deployment integration.


\noindent\textbf{Self-Reflection.} Self-reflection enables LLM agents to autonomously critique past executions and iteratively refine subsequent strategies. Reflexion \citep{shinn2023reflexion} established the foundation for this paradigm by introducing verbal reinforcement to store linguistic feedback to prevent recurring errors. Self-Refine \citep{madaan2023self} employs an iterative loop where the model generates, critiques, and corrects its own output without external supervision. Recent advancements explore more granular reflection strategies: AR \citep{wang2024devil} focuses on real-time adaptation by preparing backup moves and switching to alternatives upon primary execution failure. MIRROR \citep{guo2025mirror} further extends this by integrating inter-agent communication with intra-agent self-critique. SAMULE \citep{ge2025samule} shifts from prompting to specialized training, fine-tuning a reflection LLM to provide high-quality feedback to an executor agent. 
% We provide an analysis of similar baselines in Appendix~\ref{appendix:baseline}.